\documentclass[12pt]{article}
\usepackage{fullpage}
\usepackage{amsmath,amssymb,amsthm}

\newtheorem{lemma}{Lemma}
\newtheorem{proposition}{Proposition}
\newcommand{\R}{\mathbb R}
\newcommand{\E}{\mathbb E}
\newcommand{\1}{\mathbf 1}
\newcommand{\Var}{\operatorname{Var}}
\newcommand{\TV}{\operatorname{TV}}
\newcommand{\Law}{\operatorname{Law}}
\newcommand{\calF}{\mathcal F}

\begin{document}

\begin{center}
{\bf Uniqueness of the invariant measure}
\end{center}

\section*{Problem statement and interpretation}

Let \(E=C_0([0,1])\), \(H=L^2(0,1)\), and let \(S_t\) be the Dirichlet heat
semigroup for \(\frac12\Delta\), with kernel \(p_t(x,y)\).  I use the
Dirichlet extension of the Athreya--Butkovsky--L\^e--Mytnik (ABLM) solution
stipulated in the problem in the following standard mild form.  For the
regularized equations with \(b_n=G_{1/n}\),
\[
  u_t^n=S_t x+V_t+K_t^n,\qquad
  K_t^n=\int_0^t S_{t-r}b_n(u_r^n)\,dr,                    \tag{1}
\]
and \(u^n\Rightarrow u\) in the ABLM approximation sense, so that
\[
        P_t^n\Phi(x)\longrightarrow P_t\Phi(x),
        \qquad x\in E,\quad t>0,\quad \Phi\in C_b(E).        \tag{2}
\]
For the limit solution we write \(u_t=S_tx+V_t+K_t\), where \(K\) is the
limit drift additive functional.  The pointwise \(V(3/4)\) estimate used by
ABLM is delicate at a Dirichlet boundary.  We shall only need, and prove
below, the weaker Hilbert-space estimate
\[
 \|K_t-S_{t-s}K_s\|_{L^2(\Omega;H)}\le C_T |t-s|^{3/4},
        \qquad 0\le s<t\le T .                              \tag{3}
\]

\begin{lemma}[Hilbert drift estimate]\label{lem:Hestimate}
For every \(T<\infty\), (3) holds.  The same bound holds for \(K^n\), with a
constant independent of \(n\).
\end{lemma}

\begin{proof}
It is enough to prove the uniform estimate for \(K^n\); the limit follows by
lower semicontinuity along the ABLM approximation of the mild decomposition.

Put \(d(y)=y\wedge(1-y)\) and
\[
        \rho_\tau(y)=\int_0^\tau p_{2a}(y,y)\,da .
\]
By the reflection formula for the Dirichlet kernel,
\[
        c(\sqrt{\tau}\wedge d(y))\le \rho_\tau(y)
        \le C(\sqrt{\tau}\wedge d(y)),\qquad 0<\tau\le T.     \tag{4}
\]
Consequently
\[
\int_0^1\rho_\tau(y)^{-1/2}dy\le C\tau^{-1/4},\qquad
\sup_z\int_0^1 p_a(z,y)\rho_\tau(y)^{-1/2}dy
     \le C(a^{-1/4}+\tau^{-1/4}).                            \tag{5}
\]
The estimate needed for the sewing defect is not a false pointwise bridge
inequality, but the following \(H\)-operator bound.  Since
\(\rho_\tau^{-1}\le C(\tau^{-1/2}+d^{-1})\), Hardy's inequality on
\(H_0^1(0,1)\) and the heat-semigroup gradient estimate give
\[
        \|\rho_\tau^{-1}S_\tau f\|_H
        \le C\tau^{-1/2}\|f\|_H,\qquad f\in H .              \tag{6}
\]

For \(0\le s<t\le T\) define
\[
D^n_{s,t}=K_t^n-S_{t-s}K_s^n
\]
and the frozen one-step approximation
\[
 A^n_{s,t}=\int_s^t S_{t-r}
 b_n\!\left(S_{r-s}u_s^n+V_r-S_{r-s}V_s\right)\,dr .          \tag{7}
\]
For fixed \(n\), \(D^n\) is the semigroup-sewn integral of \(A^n\); this is
immediate by taking Riemann sums and using the smoothness of \(b_n\).

First, uniformly in \(n\),
\[
        \|A^n_{s,t}\|_{L^2(\Omega;H)}\le C|t-s|^{3/4}.        \tag{8}
\]
Indeed, condition on \(\calF_s\).  If \(s<r<q\), the two successive Gaussian
increments have conditional variances \(\rho_{r-s}(y)\) and
\(\rho_{q-r}(z)\).  Since \(b_n\) is a probability density,
\[
 \E_s\, b_n(Z_r(y))b_n(Z_q(z))
 \le C\,\rho_{r-s}(y)^{-1/2}\rho_{q-r}(z)^{-1/2},
\]
where \(Z_r=S_{r-s}u_s^n+V_r-S_{r-s}V_s\).  Using
\(\int p_a(x,y)p_b(x,z)\,dx=p_{a+b}(y,z)\), (5), and then scaling the
two-dimensional time integral over \(s<r<q<t\), gives (8):
\[
\begin{aligned}
 \E_s\|A^n_{s,t}\|_H^2
 &\le C\!\int_s^t\!\!\int_r^t\!\!
      \iint p_{2t-r-q}(y,z)
      \rho_{r-s}(y)^{-1/2}\rho_{q-r}(z)^{-1/2}\,dy\,dz\,dq\,dr  \\
 &\le C\!\int_s^t\!\!\int_r^t
      \bigl((2t-r-q)^{-1/4}+(r-s)^{-1/4}\bigr)(q-r)^{-1/4}
      \,dq\,dr
 \le C|t-s|^{3/2}.
\end{aligned}
\]

Second, if \(s<u<t\), then with the semigroup coboundary
\[
\delta A^n_{s,u,t}=A^n_{s,t}-S_{t-u}A^n_{s,u}-A^n_{u,t},
\]
we have the pathwise conditional estimate
\[
        \|\E_u\delta A^n_{s,u,t}\|_H
        \le C(t-u)^{1/2}\|D^n_{s,u}\|_H .                   \tag{9}
\]
To see this, note that for \(r>u\) the two frozen arguments differ exactly by
\(S_{r-u}D^n_{s,u}\).  Conditional on \(\calF_u\), the remaining increment is
Gaussian with variance \(\rho_{r-u}(y)\).  Therefore the derivative bound for
a one-dimensional Gaussian convolution gives, uniformly in \(n\),
\[
 \left|\E_u\{b_n(Z_r(y))-b_n(Z_r(y)+h(y))\}\right|
      \le C\rho_{r-u}(y)^{-1}h(y),\quad h=S_{r-u}D^n_{s,u}.
\]
After applying \(S_{t-r}\), (6) yields (9).

Now apply the Hilbert-valued stochastic sewing lemma of L\^e in its
semigroup form.  Let
\[
Y_n(h)=\sup_{0<t-s\le h}
       \frac{\|D^n_{s,t}\|_{L^2(H)}}{|t-s|^{3/4}} .
\]
On a fixed interval of length at most \(h\), (9) and the definition of
\(Y_n(h)\) imply
\[
 \|\E_a\delta A^n_{a,b,c}\|_{L^2(H)}
 \le C\,Y_n(h)(b-a)^{3/4}(c-b)^{1/2}
 \le C\,Y_n(h)(c-a)^{5/4}.
\]
Also \(\|\delta A^n_{a,b,c}-\E_a\delta A^n_{a,b,c}\|_{L^2(H)}
 \le C(c-a)^{3/4}\), by (8).  The sewing lemma therefore yields
\[
 \|D^n_{s,t}-A^n_{s,t}\|_{L^2(H)}
 \le C|t-s|^{3/4}+C\,Y_n(h)|t-s|^{5/4}.
\]
Together with (8) this gives \(Y_n(h)\le C+C h^{1/2}Y_n(h)\).  For \(h\)
small this absorbs the last term, uniformly in \(n\).  Larger intervals are
covered by the semigroup additivity
\(D^n_{s,t}=S_{t-r}D^n_{s,r}+D^n_{r,t}\) and the preceding local estimate
(together with \(t-s\ge h\)).  This proves (3) for \(K^n\), uniformly in
\(n\), and hence for \(K\).
\end{proof}

\section*{The reversible Gibbs measure}

Let \(\gamma\) be Brownian bridge measure on \(E\), the centered Gaussian
measure with covariance \((-\Delta_D)^{-1}\).  Define
\[
  \Lambda(v)=\int_0^1\1_{\{v(r)>0\}}\,dr,\qquad
  \pi(dv)=Z^{-1}e^{2\Lambda(v)}\,\gamma(dv).                 \tag{10}
\]
Then \(\pi\sim\gamma\) and has full support.  For \(B_n'=b_n\), with
\(B_n(a)=\int_{-\infty}^a b_n(r)\,dr\), put
\[
   \pi_n(dv)=Z_n^{-1}\exp\left(2\int_0^1B_n(v(r))\,dr\right)\gamma(dv).
\]
Since \(0\le B_n\le1\), \(B_n(a)\to\1_{\{a>0\}}\) for \(a\ne0\), and a
Brownian bridge spends zero Lebesgue time at zero, dominated convergence gives
\[
        \|\pi_n-\pi\|_{\TV}\to0 .                            \tag{11}
\]

For fixed \(n\), the mild SPDE with drift \(b_n\) is the standard gradient
diffusion associated with the closed form
\[
        {\cal E}_n(F,G)=\frac12\int_E\langle DF,DG\rangle_H\,d\pi_n .
\]
On smooth cylinder functions,
\[
 L_nF(v)=\frac12{\rm Tr}\,D^2F(v)
       +\left\langle \frac12\Delta_Dv+b_n(v),DF(v)\right\rangle_H
\]
and Gaussian integration by parts gives
\[
 \int_E G\,L_nF\,d\pi_n
 =-\frac12\int_E\langle DF,DG\rangle_H\,d\pi_n
 =\int_E F\,L_nG\,d\pi_n .                                  \tag{12}
\]
Thus \(P_t^n\) is \(\pi_n\)-invariant and reversible.

\begin{proposition}\label{prop:gap}
The measure \(\pi\) is \(P_t\)-invariant.  Moreover \(P_t\) is symmetric on
\(L^2(\pi)\), and there exists \(c>0\) such that
\[
       \Var_\pi(P_t f)\le e^{-2ct}\Var_\pi(f),
       \qquad f\in L^2(\pi),\ t\ge0 .                       \tag{13}
\]
\end{proposition}

\begin{proof}
For \(f,g\in C_b(E)\), (2), (11), and the invariance and symmetry of
\((P_t^n,\pi_n)\) imply
\[
 \int P_tf\,d\pi=\lim_n\int P_t^nf\,d\pi_n=\int f\,d\pi
\]
and
\[
 \int gP_tf\,d\pi=\lim_n\int gP_t^nf\,d\pi_n
       =\lim_n\int fP_t^ng\,d\pi_n=\int fP_tg\,d\pi .
\]
This gives invariance and symmetry, first on \(C_b(E)\), then on \(L^2(\pi)\)
by density and Jensen's inequality.

The Gaussian Poincare inequality for \(\gamma\), together with the uniform
bounds on \(d\pi_n/d\gamma\), gives by Holley--Stroock
\[
      \Var_{\pi_n}(F)\le C\,{\cal E}_n(F,F)
\]
with \(C\) independent of \(n\).  Since \(P_t^n\) is the symmetric semigroup
generated by \({\cal E}_n\),
\[
      \Var_{\pi_n}(P_t^n f)\le e^{-2t/C}\Var_{\pi_n}(f).
\]
Passing to the limit with (2) and (11), and then using density of \(C_b(E)\)
in \(L^2(\pi)\), proves (13).
\end{proof}

\section*{Fixed-time absolute continuity}

\begin{lemma}\label{lem:control}
For every deterministic \(x\in E\) and \(t>0\),
\[
        P_t(x,\cdot)\ll \gamma ,
\]
hence \(P_t(x,\cdot)\ll\pi\).
\end{lemma}

\begin{proof}
Write \(u_t=S_tx+V_t+K_t\).  Let \(r_k=t(1-2^{-k})\) and
\(\Delta_k=r_{k+1}-r_k\).  Set
\[
        D_k=K_{r_{k+1}}-S_{\Delta_k}K_{r_k}.
\]
By Lemma \ref{lem:Hestimate},
\[
        \E\|D_k\|_H^2\le C_t\,\Delta_k^{3/2}.                \tag{14}
\]
Define the predictable \(H\)-valued process
\[
 F_s=\sum_{k\ge0}\frac{\1_{(r_{k+1},r_{k+2}]}(s)}{\Delta_{k+1}}\,
        S_{s-r_{k+1}}D_k .                                  \tag{15}
\]
Then
\[
 \E\int_0^t\|F_s\|_H^2ds
 \le C_t\sum_{k\ge0}\frac{\Delta_k^{3/2}}{\Delta_{k+1}}<\infty. \tag{16}
\]
Moreover the terminal effect of this adapted control is exactly \(K_t\):
\[
 \int_0^t S_{t-s}F_s\,ds
 =\lim_{N\to\infty}\sum_{k=0}^N S_{t-r_{k+1}}D_k
 =\lim_{N\to\infty}S_{t-r_{N+1}}K_{r_{N+1}}=K_t              \tag{17}
\]
in \(H\).  Thus
\[
        u_t=S_tx+V_t+\int_0^tS_{t-s}F_s\,ds .                \tag{18}
\]

Let \(\tau_N=\inf\{s:\int_0^s\|F_r\|_H^2dr>N\}\) and
\(F^N=F\1_{[0,\tau_N]}\).  Novikov's condition holds for \(F^N\), and the
Hilbert-space Girsanov theorem gives absolute continuity of the terminal
\(H\)-law of
\[
 S_\cdot x+V_\cdot+\int_0^\cdot S_{\cdot-r}F^N_r\,dr
\]
with respect to the terminal Ornstein--Uhlenbeck law.  On \(\{\tau_N>t\}\)
this terminal value equals \(u_t\), while \(\tau_N\uparrow\infty\) a.s.;
hence \(\Law_H(u_t)\ll\Law_H(S_tx+V_t)\).

In the eigenbasis \(e_k(r)=\sqrt2\sin(k\pi r)\), the latter covariance has
eigenvalues
\[
        q_k(t)=\frac{1-e^{-k^2\pi^2t}}{k^2\pi^2},
\]
whereas \(\gamma\) has eigenvalues \((k^2\pi^2)^{-1}\).  The covariance ratios
differ from \(1\) by an exponentially square-summable sequence, and
\(S_tx\in H_0^1(0,1)\), the Cameron--Martin space of \(\gamma\).  The
Feldman--Hajek theorem gives \(\Law_H(S_tx+V_t)\sim\gamma\).  Since both
measures are supported on \(E\), the same absolute continuity holds for the
\(E\)-valued laws.
\end{proof}

\section*{Conclusion: uniqueness}

Let \(\eta\) be any invariant probability measure.  If \(\pi(A)=0\), then
Lemma \ref{lem:control} gives \(P_t(y,A)=0\) for every \(y\in E\), so
\[
        \eta(A)=\int_E P_t(y,A)\,\eta(dy)=0.
\]
Thus \(\eta\ll\pi\); write \(d\eta=h\,d\pi\).

The detailed-balance identity extends from bounded continuous functions to all
bounded Borel functions by a monotone-class argument applied to the finite
measure \(\pi(dx)P_t(x,dy)\) on the Polish space \(E\times E\).  Hence
\(P_t\) is self-adjoint also in the \(L^1\)-\(L^\infty\) duality.  For bounded
Borel \(f\),
\[
 \int f h\,d\pi=\int P_tf\,d\eta
      =\int P_tf\,h\,d\pi=\int f\,P_t h\,d\pi ,
\]
so \(P_t h=h\) in \(L^1(\pi)\).

For \(a>0\), let \(h_a=h\wedge a\).  Since \(r\mapsto r\wedge a\) is concave,
\[
        P_t h_a\le (P_t h)\wedge a=h_a .
\]
The two sides have the same \(\pi\)-integral, hence \(P_t h_a=h_a\).  Now
\(h_a\in L^2(\pi)\), and the spectral gap (13) gives
\(\Var_\pi(h_a)=0\).  Therefore \(h_a\) is constant for every \(a\).  Letting
\(a\uparrow\infty\) and using \(\int h\,d\pi=1\), we obtain \(h\equiv1\).
Thus \(\eta=\pi\).

Consequently \(P_t\) has at most one invariant probability measure; in fact the
unique invariant measure is the Gibbs measure \(\pi\) in (10).

\begin{thebibliography}{9}

\bibitem{ABLM}
S. Athreya, O. Butkovsky, K. L\^e and L. Mytnik,
\emph{Well-posedness of stochastic heat equation with distributional drift
and skew stochastic heat equation}, Comm. Pure Appl. Math. 77 (2024),
2708--2777.

\bibitem{Bogachev}
V. I. Bogachev,
\emph{Gaussian Measures}, American Mathematical Society, 1998.

\bibitem{DPZ}
G. Da Prato and J. Zabczyk,
\emph{Ergodicity for Infinite Dimensional Systems}, Cambridge University
Press, 1996.

\bibitem{BZ}
S. K. Bounebache and L. Zambotti,
\emph{A skew stochastic heat equation}, J. Theoret. Probab. 27 (2014),
168--201.

\end{thebibliography}

\end{document}
